\clearpage
\section{\edtext{Four applications under test}{Kiểm tra bốn ứng dụng}}\label{sec:study}
\edtext{We ask when a guarantee about an action survives the measurement resolution, communication schedule or observation horizon of a particular system. Two experiments select a window in a sequence, one reconstructs a moving integer, and one tracks the output of a register. We report synthetic workloads, an observational replay and finite implementation tests separately because each answers a different question.}
{Phần này xét khi nào một quyết định còn được bảo đảm với độ phân giải của phép đo, lịch truyền tin hoặc số bước quan sát cho trước. Hai thí nghiệm chọn cửa sổ trên một dãy, một thí nghiệm tái dựng một đại lượng nguyên biến đổi và một thí nghiệm theo dõi đầu ra của thanh ghi. Dữ liệu tải tổng hợp, bản ghi từ hệ vận hành và phép kiểm hữu hạn trên mô hình thực thi cung cấp ba loại bằng chứng khác nhau.}

\subsection{\edtext{What a window certificate measures}{Chứng nhận cho một cửa sổ đo điều gì?}}
\edtext{Partition each trace into batches and consider all windows of a fixed width $w$ within each batch. Select the earliest maximizer $I$ of the rounded sum $\widehat S$. Keep every tied maximizer when assessing correctness against the reference values. If the rounding step is $q$, each coordinate differs from its reported value by at most $q/2$. For an allowable loss $\delta\ge0$ in the sum, define}
{Chia mỗi bản ghi thành các nhóm, rồi xét mọi cửa sổ có cùng độ dài $w$ trong một nhóm. Chọn cửa sổ $I$ bắt đầu sớm nhất trong số các cửa sổ đạt tổng lớn nhất theo số đã làm tròn $\widehat S$. Khi đối chiếu với dữ liệu tham chiếu, mọi cửa sổ đồng tối ưu đều được tính đúng. Với bước làm tròn $q$, sai số mỗi tọa độ không quá $q/2$. Cho phép hụt tối đa $\delta\ge0$ trên tổng và đặt}
\[
 M(I)=\min_{J\ne I}\left[\widehat S(I)-\widehat S(J)
                  -\tfrac q2 |I\mathbin{\triangle}J|\right].
\]
\edtext{The certificate accepts if $M(I)\ge-\delta$. Indeed, in the difference of two sums each shared coordinate cancels. Every remaining error contributes at most $q/2$ against $I$. Consequently $\max_J S(J)-S(I)\le\max(0,-M(I))$, by the argument of Theorem~\ref{thm:heterogeneous}. The baseline bounds each sum separately and subtracts $wq$ for every competitor. The two methods select the same nominal window. Their acceptance counts measure how often each can justify that selection.}
{Chấp nhận nếu $M(I)\ge-\delta$. Trong hiệu hai tổng, sai số ở từng tọa độ chung triệt tiêu. Từng tọa độ còn lại có thể làm $I$ bất lợi thêm nhiều nhất $q/2$. Do đó $\max_J S(J)-S(I)\le\max(0,-M(I))$, theo lập luận của Định lý~\ref{thm:heterogeneous}. Cách đối chứng chặn sai số riêng cho từng tổng nên phải trừ $wq$ với mọi đối thủ. Hai cách cùng chọn một cửa sổ từ số đo danh nghĩa. Số lần chấp nhận cho biết mỗi cách bảo đảm được lựa chọn đó bao nhiêu lần.}

\subsection{\edtext{Server latency and a held-out CDN replay}{Độ trễ máy chủ và phép thử trên bản ghi CDN giữ riêng}}
\edtext{In SimPy, a single FIFO server serves 48 requests per run with an elevated service demand during a burst. The action identifies six consecutive requests with the greatest total latency. Across 100 seeded runs, increasing the rounding step from 2 to 8 to 24 ms reduces exact certificates from 56 to 9 to 1. Always choosing the nominal winner gives 7, 46 and 83 strictly suboptimal decisions, respectively. The separate-sum baseline certifies none in this simulation. The data are generated by a single queue~\cite{study-simpy}.}
{Trong SimPy, một máy chủ phục vụ 48 yêu cầu mỗi lượt theo thứ tự đến; thời gian phục vụ tăng trong một đợt tải. Nhiệm vụ là tìm sáu yêu cầu liên tiếp có tổng độ trễ lớn nhất. Qua 100 lượt với hạt giống xác định, khi bước làm tròn tăng từ 2 lên 8 rồi 24 ms, số lựa chọn được chứng nhận chính xác giảm từ 56 xuống 9 rồi 1. Nếu luôn chốt theo số làm tròn thì số lựa chọn hụt đáp án lần lượt là 7, 46 và 83. Cách chặn riêng từng tổng không chứng nhận được trường hợp nào trong mô phỏng này. Dữ liệu ở đây do một hàng đợi sinh ra~\cite{study-simpy}.}

\edtext{We next used published observations of time to first byte (TTFB) on a Wikimedia text cache~\cite{study-wikimedia}. The reference is the value in the released trace. A pilot on days 00 and 14 used 19,200 requests. At $q=8$ ms the overlap-aware rule certified 204 of 400 groups, compared with 186 for the baseline. The nominal choice was suboptimal in 103 groups. Before reading the other files, we fixed the days (02, 09 and 18), hours (00, 08 and 16 relative to the start of each published day), window width, rounding, ties, metrics and case-selection rules. We took the first 9,600 consecutive records from each cutoff. Each slice spans 54--200 seconds. Within the same published second we retain file order. This yields 1,800 nonoverlapping groups of 48, each with 43 candidate windows. Adjacent groups can share operating conditions.}
{Tiếp đó, chúng tôi dùng số đo thời gian tới byte đầu tiên (TTFB) do Wikimedia công bố cho bộ nhớ đệm phục vụ trang văn bản~\cite{study-wikimedia}. Giá trị tham chiếu là số đo trong bản ghi công bố. Thử thăm dò gồm 19.200 yêu cầu ở ngày 00 và 14. Với $q=8$ ms, cách xét phần giao chứng nhận đúng 204 trong 400 nhóm, so với 186 của cách đối chứng. Lựa chọn theo số làm tròn hụt đáp án ở 103 nhóm. Trước khi đọc các tệp còn lại, chúng tôi ấn định ngày 02, 09, 18; các mốc giờ 00, 08, 16 tính từ đầu mỗi ngày trong bộ dữ liệu; độ dài cửa sổ, cách làm tròn, xử lý hòa, chỉ tiêu và quy tắc chọn ca. Lấy 9.600 bản ghi liên tiếp đầu tiên kể từ mỗi mốc. Mỗi lát kéo dài 54--200 giây. Trong cùng một giây, giữ thứ tự của tệp nguồn. Kết quả gồm 1.800 nhóm không chồng nhau, mỗi nhóm 48 yêu cầu và 43 cửa sổ ứng viên. Hai nhóm liền nhau có thể cùng chịu một điều kiện vận hành.}

\begin{center}
\begin{tabular}{rrrrr}
\toprule
$\delta$ (ms) & \edtext{Certified}{Chứng nhận} & \edtext{Baseline}{Đối chứng} &
\edtext{Loss $>\delta$}{Hụt $>\delta$} & \edtext{Certified, loss $>\delta$}{Đã nhận, hụt $>\delta$}\\
\midrule
0 &855&766&486&0\\
1 &855&766&10&0\\
8 &1339&797&1&0\\
\bottomrule
\end{tabular}
\end{center}
\edtext{All counts refer to the same 1,800 held-out groups (86,400 requests). Loss is the shortfall in the sum of six TTFBs in milliseconds. The fourth column evaluates a policy that always selects the nominal winner. At zero tolerance the tighter rule adds 89 certificates. At 1 ms only ten nominal choices exceed tolerance, yet the certificate count is unchanged: rounded scores and pairwise error allowances lie on an 8 ms grid. The positive thresholds probe sensitivity; neither was set by an operator.}
{Mẫu số của mỗi dòng là cùng 1.800 nhóm giữ riêng (86.400 yêu cầu). Mức hụt là chênh lệch trên tổng TTFB của sáu yêu cầu, tính bằng mili giây. Cột thứ tư chấm cách luôn chọn cửa sổ danh nghĩa. Ở mức hụt bằng không, cách xét phần giao chứng nhận thêm 89 nhóm. Cho phép hụt 1 ms thì chỉ còn mười lựa chọn vượt ngưỡng, nhưng số chứng nhận giữ nguyên vì tổng sau làm tròn và mức bù sai số giữa hai cửa sổ đều nằm trên lưới bước 8 ms. Hai ngưỡng dương dùng để khảo sát độ nhạy; người vận hành chưa xác lập ngưỡng nghiệp vụ.}

\edtext{Among the 486 suboptimal choices, 341 shared five of six requests with an optimal window and six shared none with any optimum. The median shortfall among mistakes was 0.0355 ms. The largest, 12.089 ms, occurred for disjoint windows. Nine alternating pairs of timed Python passes over the held-out batches gave a median additional 8.38 microseconds per batch for certification, and a repeat gave 9.92 microseconds. Both timings start from rounded inputs and exclude I/O and comparison with the reference. The supplement contains eight blinded A/B cases, including controls, for assessing whether a changed location changes triage. No operator judgments have been collected. The available evidence establishes positional and score differences; incident identity remains unobserved.}
{Trong 486 lựa chọn hụt đáp án, 341 cửa sổ trùng năm trong sáu yêu cầu với một cửa sổ tối ưu, còn sáu cửa sổ rời mọi đáp án tối ưu. Trung vị mức hụt trong các lần chọn sai là 0,0355 ms. Trường hợp hụt nhiều nhất, 12,089 ms, có hai cửa sổ rời nhau. Chín cặp chạy Python đổi thứ tự luân phiên cho thấy thời gian tăng thêm có trung vị 8,38 micro giây mỗi nhóm để tính chứng nhận; chạy lặp lại được 9,92 micro giây. Cả hai phép đo bắt đầu từ số đã làm tròn và loại thời gian đọc ghi cũng như đối chiếu đáp án. Phụ lục dữ liệu có tám ca A/B ẩn cách gán nhãn, gồm cả ca đối chứng, để đánh giá việc đổi cửa sổ có đổi hướng điều tra hay không. Chưa thu nhận ý kiến người trực hệ thống. Số liệu cho biết hai cửa sổ khác vị trí và khác điểm bao nhiêu; danh tính của sự cố vẫn chưa được quan sát.}

\subsection{\edtext{Power-system windows}{Chọn cửa sổ trong hệ thống điện}}
\edtext{Using pandapower and its 33-bus feeder case, we solved AC power flow for 24 synthetic demand profiles, each with sixteen 15-minute slots. All 384 solves converged. Substation import, including network losses, was integrated over each slot under constant power and rounded to Wh for the reference. The objective was the largest total import in two consecutive slots~\cite{study-pandapower}.}
{Dùng pandapower và mô hình lưới 33 nút đi kèm, chúng tôi giải trào lưu công suất xoay chiều cho 24 chuỗi tải tổng hợp; mỗi chuỗi có 16 khoảng dài 15 phút. Cả 384 lần giải đều hội tụ. Công suất nhập tại trạm, gồm cả tổn thất trên lưới, được coi là cố định trong mỗi khoảng để tính điện năng và làm tròn tới Wh làm tham chiếu. Mục tiêu là tìm hai khoảng liên tiếp có tổng điện năng nhập lớn nhất~\cite{study-pandapower}.}

\begin{center}\begin{tabular}{rrrr}
\toprule $q$ (kWh) & \edtext{Certified}{Chứng nhận} & \edtext{Baseline}{Đối chứng} & \edtext{Nominal errors}{Lựa chọn sai}\\
\midrule 10&22&18&1\\50&11&1&6\\150&8&0&9\\\bottomrule
\end{tabular}\end{center}
\edtext{Each row represents the same 24 profiles. Every accepted window was optimal on its reference trace. At 50 kWh the overlap-aware rule accepted ten additional profiles. In one rejected case the nominal window started one slot early and its import was 33.328 kWh below the optimum. These are retrospective window decisions. A battery dispatch policy would additionally need forecasts available at the decision time, state of charge, conversion losses and power limits.}
{Mỗi dòng dùng cùng 24 chuỗi tải. Mọi cửa sổ được chứng nhận đều tối ưu theo chuỗi tham chiếu. Ở bước 50 kWh, cách xét phần giao chứng nhận thêm mười chuỗi. Trong một ca bị từ chối chứng nhận, cửa sổ danh nghĩa bắt đầu sớm một khoảng và tổng điện năng nhập thấp hơn đáp án 33,328 kWh. Đây là lựa chọn hồi cứu. Muốn điều độ pin, còn cần dự báo sẵn có lúc ra quyết định, trạng thái sạc, tổn hao chuyển đổi và giới hạn công suất.}

\subsection{\edtext{Sensor communication with a deadline}{Truyền tin cảm biến dưới hạn chót}}
\edtext{Two FIFO ports in ns.py implement the integer model of Theorem~\ref{B-thm:integer} with $n=8$, $b=7$, $d=4$ and two early labels when an early message is sent. We enumerate all 576 admissible triples of old state, receiver information and fresh state for each schedule and load condition. Compare one source bit before and one after the fresh measurement with two bits after it. Receiver feedback is sent at 10 ms, the fresh sample is taken at 20 ms, and the deadline is 26 ms. Each packet carries a four-byte header; payload is rounded up to whole bytes. Each direction serves eight bits per millisecond with zero propagation delay~\cite{study-ns}.}
{Hai cổng FIFO trong ns.py thực thi mô hình nguyên của Định lý~\ref{B-thm:integer} với $n=8$, $b=7$, $d=4$ và hai nhãn sớm khi nguồn gửi trước. Vét cạn 576 bộ ba hợp lệ gồm trạng thái cũ, thông tin bên nhận và trạng thái mới cho mỗi tổ hợp lịch gửi và điều kiện tải. So sánh một bit nguồn trước và một bit sau phép đo mới với hai bit nguồn gửi sau phép đo. Bên nhận phản hồi ở 10 ms, mẫu mới được lấy ở 20 ms và hạn chót là 26 ms. Mỗi gói có phần đầu bốn byte, phần dữ liệu được làm tròn lên byte. Mỗi chiều truyền tám bit trên mili giây và giả định thời gian lan truyền bằng không~\cite{study-ns}.}

\begin{center}\begin{tabular}{lrrrr}
\toprule \edtext{Source bits}{Bit nguồn} & \edtext{Max.\ error}{Sai số cực đại} & \edtext{Total bytes}{Tổng byte} & \edtext{Idle (ms)}{Rảnh (ms)} & \edtext{Busy (ms)}{Bận (ms)}\\
\midrule $1+1$&$5/2$&15&25&34\\ $0+2$&$3/2$&10&25&34\\\bottomrule
\end{tabular}\end{center}
\edtext{Byte counts include both directions of the protocol. Adding a ten-byte background packet at 19 ms changes completion from 25 to 34 ms in both schedules. Every triple then misses the deadline, although the reconstruction error with respect to the sample at 20 ms is unchanged. A guarantee for the state at receipt would require an additional drift bound.}
{Tổng byte tính cả hai chiều giao thức. Khi đưa một gói nền mười byte vào hàng đợi ở 19 ms, thời điểm hoàn tất của cả hai lịch tăng từ 25 lên 34 ms. Khi ấy mọi bộ ba đều trượt hạn, dù sai số tái dựng đối với mẫu ở 20 ms vẫn giữ nguyên. Muốn bảo đảm cho trạng thái lúc nhận kết quả, cần thêm cận biến động sau khi lấy mẫu.}

\subsection{\edtext{Finite-horizon hardware monitoring}{Giám sát phần cứng với số bước hữu hạn}}
\edtext{PyRTL simulates a 16-bit register whose output is its low bit. At each step the circuit rotates right or XORs with the known mask \texttt{0x8001}. We compare its output with a monitor initialized using only the four low bits. For each of 160 initial states, all eight length-three command words are exercised. All 5,120 output comparisons agree (1,280 trajectories, including the initial observation). Three successive rotations distinguish all 16 low-nibble patterns, giving a four-bit lower bound. At four rotations, states 0 and 16 yield respective output sequences $00000$ and $00001$. Thus extending the certified horizon by one step exposes information discarded at initialization. The trial measures functional behavior; area, power and timing on a fabricated device remain open~\cite{study-pyrtl}.}
{PyRTL mô phỏng thanh ghi 16 bit có đầu ra là bit thấp nhất. Mỗi bước, mạch xoay phải hoặc XOR với mặt nạ đã biết \texttt{0x8001}. So đầu ra với bộ theo dõi được khởi tạo bằng bốn bit thấp. Với mỗi trạng thái trong 160 trạng thái đầu, chạy đủ tám chuỗi ba lệnh. Cả 5.120 lần so đầu ra đều trùng nhau (1.280 quỹ đạo, tính cả quan sát ban đầu). Ba lần xoay liên tiếp phân biệt đủ 16 mẫu bốn bit thấp, cho cận dưới bốn bit. Sau bốn lần xoay, hai trạng thái 0 và 16 lần lượt cho chuỗi đầu ra $00000$ và $00001$. Mở rộng thời hạn bảo đảm chỉ một bước đã bộc lộ thông tin bị bỏ lúc khởi tạo. Phép thử đo hành vi chức năng; diện tích mạch, điện năng và thời gian thực thi trên chip chưa được đo~\cite{study-pyrtl}.}

\edtext{The computational artifact in \texttt{supplement/study} contains the methods, inputs, machine-readable outputs, software versions and replay instructions. Separate repeated runs reproduced the deterministic outcomes. The four endpoints are different quantities: certificate coverage, import shortfall, reconstruction error at a sampling time, and agreement of output bits. Each conclusion applies to its specified endpoint.}
{Bộ tái lập trong \texttt{supplement/study} cung cấp phương pháp, đầu vào, đầu ra đọc được bằng máy, phiên bản phần mềm và lệnh chạy lại. Các lượt chạy độc lập tái tạo cùng kết quả tất định. Bốn chỉ tiêu có ý nghĩa riêng: tỷ lệ quyết định đủ chứng nhận, mức hụt điện năng nhập, sai số tái dựng ở thời điểm lấy mẫu và sự trùng khớp của các bit đầu ra. Mỗi kết luận gắn với chỉ tiêu đã nêu.}
